\section{Guía introductoria: cómo juzgar si un sistema de AI for Math es confiable}

Al estudiar este episodio, la pregunta «¿es confiable un sistema de AI for Math?» puede dividirse en seis preguntas. Primera: ¿puede comprender con precisión los objetos matemáticos y los enunciados de los teoremas, en lugar de limitarse a hacer una paraphrase en lenguaje natural? Segunda: ¿puede convertir una demostración informal en una forma que un proof assistant como Lean/Coq pueda verificar? Tercera: cuando una demostración falla, ¿puede leer el goal state y repararla, en lugar de reescribir un texto alucinado? Cuarta: ¿puede aprovechar una biblioteca de teoremas, en lugar de empezar desde cero cada vez? Quinta: ¿puede acelerar la investigación de los matemáticos, y no solo resolver problemas ya conocidos? Sexta: ¿puede explicar por qué una demostración es significativa?

\begin{importantbox}{El método de las seis preguntas}
Si un sistema de AI for Math solo cumple las dos primeras preguntas, se parece más a una herramienta de formalization; si cumple la tercera y la cuarta, empieza a tener capacidades de Agent; solo si cumple la quinta y la sexta se acerca a ser un colaborador de investigación matemática.
\end{importantbox}

\subsection{Por qué importan estas seis preguntas}

Convierten la pregunta difusa «¿sabe matemáticas la IA?» en capacidades verificables. Comprender los objetos corresponde a la capacidad semántica; la formalización corresponde a la capacidad de lenguaje y de sistema de tipos; leer el goal state corresponde a la retroalimentación del entorno; aprovechar la biblioteca de teoremas corresponde a la recuperación y la memoria; acelerar la investigación corresponde al valor de producto; explicar el significado corresponde a la sensibilidad estética matemática. Un AI for Math de alto nivel requiere que todas estas capacidades se presenten juntas.

\subsection{Resumen de la sección}

Estas seis preguntas atraviesan el resto de la nota. Aun sin conocer de antemano los detalles de Lean, el lector puede usarlas para juzgar si las distintas líneas de AI for Math están resolviendo problemas, construyendo demostraciones o intentando entrar en la investigación propiamente dicha.
